\documentclass[UTF8,a4paper,12pt]{ctexart}

\usepackage{amsmath,amssymb,bm}
\usepackage{booktabs}
\usepackage{geometry}
\usepackage{longtable}
\usepackage{array}
\usepackage{enumitem}
\usepackage{microtype}
\usepackage{hyperref}

\geometry{left=2.6cm,right=2.6cm,top=2.5cm,bottom=2.5cm}
\setlength{\parindent}{2em}
\setlength{\parskip}{0.35em}
\renewcommand{\arraystretch}{1.35}
\setlist{nosep}
\hypersetup{hidelinks}

\title{方法与实验部分}
\author{}
\date{}

\begin{document}

\maketitle

\section{3 队列感知的多线换乘站疏散路径组织模型}

大型多线换乘站的应急疏散效率受多来源客流汇聚过程控制。站台候车乘客、列车到达乘客、站厅滞留乘客和跨线换乘乘客并非沿相互独立的通道离站，而是在楼梯、扶梯、换乘通道、闸机和出口前区域等共享设施处交织。对龙阳路站这类多线换乘站而言，换乘客流虽然不是本文唯一评价对象，却是改变共享设施负荷时空分布的重要来源。其路径通常跨越不同线路的公共区和垂向连接设施，并会与站台出站客流、列车到达客流共同形成设施排队、回溢和尾部滞留。

为刻画上述过程，本文构建队列感知的动态路径组织模型 AdaptiveQueueAwareAStar，简称 AA。AA 的目标不是为车站生成一条固定疏散路线，也不是单纯比较某条路径是否更短，而是在多源客流不断进入站内网络时，预测各批次到达共享设施时的队列状态，并据此组织可由设施容量实际执行的疏散路径。模型由“路径决策层”和“设施服务执行层”组成：路径决策层负责根据未来设施负荷评价候选路径，设施服务执行层负责按照统一容量约束释放实际通行人数。两层耦合后，路径选择、设施排队和下游空间接收被放在同一疏散过程中处理。

\subsection{3.1 问题描述与模型总体框架}

设车站疏散过程在离散时刻 \(t\) 的系统状态为
\[
\Omega(t)=\left(G,B(t),Q(t),N(t),A(t)\right),
\]
其中，\(G\) 为站内疏散网络，\(B(t)\) 为进入路径选择阶段的客流批次集合，\(Q(t)\) 为共享设施队列状态，\(N(t)\) 为空间节点占用状态，\(A(t)\) 为已经承诺但尚未完全执行的未来到达事件集合。AA 在每一个时间步读取 \(\Omega(t)\)，对可决策批次依次搜索疏散路径，并将已确定路径对共享设施产生的未来到达需求写回系统状态。

模型输入包括车站拓扑、设施几何、设施服务能力、来源客流批次、候选出口、当前队列、当前空间占用和已承诺到达事件。模型输出包括每个可决策批次的候选路径、下一步移动意图、预计进入关键设施的时刻、预计等待时间以及对未来设施到达事件的更新。实际通过设施的人数不由路径搜索直接决定，而由设施服务执行层在同一时间步内统一分配。

AA 的模型框架包含四个相互衔接的部分。第一，站内疏散网络与共享设施容量模型，用于表达站台、换乘通道、楼扶梯、闸机和出口之间的连接关系，以及多条路径对同一设施服务能力的竞争。第二，多来源批次与到达负荷预测模型，用于记录不同线路和换乘关系的客流进入时刻，并预测批次到达设施时的等待状态。第三，AA 路径组织模型，用物理行走时间、设施等待、批次自身服务时间和空间回溢等待构造时间依赖路径代价。第四，疏散效果评价方法，用整体疏散效率、尾部滞留、关键设施瓶颈、来源组分解和微观仿真趋势一致性评价路径组织效果。

\subsection{3.2 站内疏散网络与共享设施容量模型}

将车站公共区抽象为有向图
\[
G=(V,E),
\]
其中，节点集合 \(V\) 表示线路站台、站厅区域、换乘连接节点、楼梯、扶梯、闸机、出口前区域和地面出口，边集合 \(E\) 表示相邻空间或设施之间允许乘客移动的方向性连接。每条边 \(e=(i,j)\) 包含长度 \(l_e\)、有效宽度 \(w_e\)、设施类型、行进方向和容量属性；每个节点 \(v\) 包含面积 \(A_v\)、当前人数 \(N_v(t)\)、节点类型和空间接收上限。

在多线换乘站中，设施容量不能按边独立使用。同一楼梯、同一扶梯、同一换乘通道或同一闸机组可能同时服务多个线路、多个站台和多个换乘方向。为描述这种共享关系，定义边到设施资源的映射
\[
\rho(e)=r,
\]
其中，\(r\) 表示边 \(e\) 需要消耗的共享设施资源。映射到同一资源 \(r\) 的多条边共同竞争该资源的服务能力；若边只表示普通空间移动而不消耗独立瓶颈设施，则其资源映射为空。该处理使换乘客流、站台客流和列车到达客流在同一设施上的叠加能够被显式记录。

设共享设施 \(r\) 的服务率为 \(\mu_r\)，单位为人每秒，仿真步长为 \(\Delta t\)。在时间步 \(t\) 内，设施 \(r\) 可释放的整数服务能力为
\[
C_r(t)=\left\lfloor \mu_r\Delta t+\xi_r(t)\right\rfloor ,
\]
其中，\(\xi_r(t)\) 为上一时间步未形成完整整数人的容量余量。设 \(A_r(t)\) 为时间步 \(t\) 内请求进入设施 \(r\) 的人数，\(Q_r(t)\) 为设施入口等待队列，则队列递推为
\[
Q_r(t+\Delta t)=\mathrm{max}\left(Q_r(t)+A_r(t)-C_r(t),0\right).
\]
当多个来源组在同一时间步到达同一资源时，未获得服务能力的人数保留在上游等待，并在后续时间步继续参与容量竞争。

设施通过能力还受到下游空间接收限制。设节点 \(v\) 的空间接收上限为 \(\bar N_v\)，当前占用人数为 \(N_v(t)\)，当前时间步已经被其他移动请求预留进入的人数为 \(B_v(t)\)，则新进入人数 \(y_v(t)\) 满足
\[
N_v(t)+B_v(t)+y_v(t)\leq \bar N_v .
\]
该约束用于描述高负荷条件下由下游通道、站厅或出口前区域引起的回溢。即使上游设施仍有名义服务能力，若下游空间不能接收，相关移动仍应停留在上游。

普通移动边的物理行走时间由边长和速度决定。当前实现采用速度-密度关系计算通行速度：
\[
v(k)=
\begin{cases}
v_f, & k\leq k_f,\\
\mathrm{max}(v_f-ak,0), & k_f<k<k_j,\\
0, & k\geq k_j .
\end{cases}
\]
其中，\(k\) 为边内客流密度，\(v_f\) 为自由流速度，\(k_f\) 为自由流密度阈值，\(a\) 为速度下降斜率，\(k_j\) 为拥挤密度上限。本文采用的 \(v_f=1.427\ \mathrm{m/s}\)、\(k_f=0.2\ \mathrm{p/m^2}\)、\(a=0.3549\) 和 \(k_j=4.0\ \mathrm{p/m^2}\) 来自对改进 A* 疏散仿真文献中速度-密度关系的复现，不作为龙阳路站实测速度。该关系用于计算物理行走时间，适用边界分析不以速度-密度阈值为核心变量。

\subsection{3.3 多来源批次与到达负荷预测模型}

设客流来源组集合为 \(\mathcal{G}\)。来源组用于区分不同线路站台、列车到达、站厅滞留和跨线换乘客流。对换乘客流，模型进一步保留其来源线路和换乘关系，使其在共享设施上的叠加可以与站台出站流、列车到达流分开统计。一个进入路径选择阶段的客流批次记为
\[
b=(g,n_b,x_b,t_b),
\]
其中，\(g\) 为来源组，\(n_b\) 为批次人数，\(x_b\) 为批次当前节点，\(t_b\) 为批次进入路径选择阶段的时刻。采用批次作为路径决策单元，可以保留列车到达、站台释放和换乘客流进入公共区的时间结构。

AA 的关键不在于读取当前设施队列，而在于预测批次真正到达共享设施时的队列状态。对共享设施 \(r\)，定义从当前时刻 \(t\) 开始已经承诺的未来到达事件集合为
\[
\mathcal{A}_r(t)=\{(\eta_j,n_j)\},
\]
其中，\(\eta_j\) 为已确定批次预计进入设施 \(r\) 的时刻，\(n_j\) 为对应人数。若候选批次 \(b\) 预计在时刻 \(\eta\) 到达设施 \(r\)，则设施到达队列估计为
\[
\widehat Q_r(\eta)=\mathcal{Q}\left(Q_r(t),\mathcal{A}_r(t),\mu_r,t,\eta\right).
\]
队列演化算子 \(\mathcal{Q}\) 在 \(t\) 到 \(\eta\) 的区间内释放服务能力，并在已登记到达事件发生时加入相应人数。若 \(\mu_r>0\)，候选批次在设施 \(r\) 处的预计等待时间为
\[
W_r(\eta)=\frac{\widehat Q_r(\eta)}{\mu_r}.
\]

未来到达事件包括两类。第一类是已经进入执行过程、到达时间可以由当前移动状态推算的批次；第二类是同一决策轮中已经完成 AA 搜索、但尚未实际到达设施的前序批次。当前序批次路径确定后，其预计进入共享设施的时间和人数立即登记为到达事件，使后续批次能够在搜索时看到这一已经承诺的设施负荷。由此，同一时间步内的批次不再基于同一份未更新状态独立选择路径，而是形成有先后顺序的设施竞争关系。

\subsection{3.4 AA 路径组织模型}

对批次 \(b\) 的候选路径 \(P=(e_1,e_2,\ldots,e_m)\)，设第 \(i\) 条边的预计进入时刻为 \(\eta_i\)，对应共享设施为 \(r_i=\rho(e_i)\)。AA 的路径代价定义为
\[
J_b(P)=\sum_{i=1}^{m}\left[\tau_{e_i}(\eta_i)+W_{r_i}(\eta_i)+S_{r_i}(n_b)+H_{e_i}(\eta_i)\right],
\]
其中，\(\tau_{e_i}\) 为物理行走时间，\(W_{r_i}\) 为到达设施时的预计排队等待，\(S_{r_i}(n_b)\) 为当前批次自身占用服务能力形成的平均服务时间，\(H_{e_i}\) 为下游空间接收不足或回溢造成的预计等待。若边不对应共享设施，则设施等待和服务时间项取 0。

路径代价按预计到达时刻递推。前一段路径上的行走和等待会改变后一段路径的到达时间，后一段设施的等待又取决于新的到达时间。因此，AA 不是静态最短路，也不是只根据当前密度重新赋权的即时路径搜索，而是面向设施到达负荷的时间依赖路径组织方法。

路径搜索标签定义为
\[
L=(v,\eta,g,\pi),
\]
其中，\(v\) 为当前节点，\(\eta\) 为预计到达当前节点的时刻，\(g\) 为累计路径代价，\(\pi\) 为前驱标签。标签从节点 \(u\) 沿边 \(e=(u,v)\) 扩展时，先计算物理行走时间，再根据预计进入设施时刻查询共享设施队列和下游空间接收状态，得到新的到达时刻与累计代价。两个标签即使到达同一节点，只要预计到达时刻不同，后续设施队列状态就可能不同，因此不能简单按节点编号合并为一个静态标签。

对仍处于选择阶段的批次，AA 允许在满足稳定性条件时重新评估路径。设保留当前路径的剩余代价为 \(C_{\mathrm{stay}}\)，候选切换路径代价为 \(C_{\mathrm{switch}}\)，相对收益为
\[
R=\frac{C_{\mathrm{stay}}-C_{\mathrm{switch}}}{C_{\mathrm{stay}}}.
\]
仅当候选路径有效、下一步移动方向不同、批次仍处于可重评估阶段、\(C_{\mathrm{switch}}<C_{\mathrm{stay}}\)，且 \(R\) 超过防摆动阈值时，路径切换才被接受。已经进入边内移动或已经获得设施服务能力的批次继续执行已承诺移动，不在途中瞬时改路。

\subsection{3.5 疏散效果评价方法}

本文从整体效率、设施瓶颈、来源组分解和跨模型一致性四个层面评价路径组织效果。整体效率指标包括 \(T_{95}\)、\(T_{100}\)、平均疏散时间、累计停滞人秒、总移动距离和计算耗时。其中，\(T_{95}\) 用于刻画系统进入尾部清空阶段前的效率，\(T_{100}\) 用于描述最后一批乘客离站时间，累计停滞人秒用于衡量乘客处于排队、阻塞和等待状态的暴露程度。

设施瓶颈指标包括关键楼扶梯、换乘通道、闸机和出口前区域的通过人数、峰值队列、队列人秒、首次排队时间、最后排队时间和利用率。来源组分解指标按线路站台、列车到达、站厅滞留和换乘来源统计清空时间、设施通过量、出口选择和队列暴露，用于说明换乘客流如何与其他来源客流共同形成共享设施负荷。跨模型一致性指标用于比较宏观网络模型与 Pathfinder 微观仿真的累计疏散曲线、拥堵位置和出口利用趋势。

\section{4 AA 模型求解与共享容量执行}

第 3 章定义了 AA 的模型组成，本章说明模型如何在离散疏散仿真中被求解和执行。这里的求解不是一次性得到全站固定路线，而是在客流批次持续进入站内网络的过程中，循环完成状态读取、路径搜索、事件登记、容量执行和指标更新。

\subsection{4.1 模型状态与搜索标签转换}

在仿真时刻 \(t\)，系统首先读取当前状态 \(\Omega(t)\)，包括节点占用、设施队列、在途批次、可决策批次和已登记的未来到达事件。对每个可决策批次，AA 将路径组织模型转化为时间依赖 A* 标签搜索。搜索入口为批次当前节点，目标为可用出口集合，边扩展过程调用同一仿真执行层中的物理行走时间、设施服务能力和空间接收状态。

标签扩展包含三类计算。第一，根据边长、设施类型和速度-密度关系计算从 \(u\) 到 \(v\) 的物理行走时间。第二，若边对应共享设施，则根据预计进入时刻查询 \(\widehat Q_r(\eta)\)，计算到达时刻设施等待和当前批次自身服务时间。第三，检查下游节点或边空间的接收条件，估计空间阻塞等待。三类成本共同更新标签的到达时刻和累计代价。

AA 的启发式函数只采用静态自由流条件下的时间下界，用于减少搜索空间；真实路径代价仍由物理行走时间、设施等待、批次服务时间和空间回溢等待共同决定。模型不缓存动态完整路径，也不使用固定 \(K\) 条候选路径。静态缓存仅用于自由流启发式下界、共享资源映射和到达事件索引，避免在预测状态已经变化后复用前序批次的动态路径。

\subsection{4.2 AA 路径搜索与事件登记}

AA 以离散时间步推进。每一时间步内，系统先更新上一时间步已经完成移动的乘客，再生成本时间步进入选择阶段的自然到达批次。可决策批次按照到达时间、来源组、当前节点和批次编号排序。排序的作用是保证同一轮内的前序批次先登记其未来设施到达事件，后续批次能够读取已经更新的预测状态。

对排序后的每个批次，AA 执行时间依赖 A* 搜索，得到候选路径、预计代价和路径上关键设施的进入时刻。若批次路径被接受，模型将其预计进入共享设施的时刻、人数、来源组和批次编号写入未来到达事件集合。该事件只表示路径组织层面的承诺需求，不代表乘客已经实际通过设施。实际通行仍需在共享容量执行层中获得服务能力。

一个时间步的 AA 路径搜索流程可概括为以下步骤：
\begin{enumerate}
    \item 读取当前节点占用、设施队列、在途批次和未来到达事件。
    \item 生成站台释放、列车到达、站厅滞留和换乘来源的自然到达批次。
    \item 对可决策批次进行确定性排序。
    \item 对当前批次执行时间依赖 A* 搜索，计算路径代价分解。
    \item 将已接受路径的共享设施到达时刻和人数登记为未来事件。
    \item 对满足稳定性条件的仍可决策批次执行路径重评估。
    \item 将所有移动意图提交给共享容量执行层。
\end{enumerate}

\subsection{4.3 共享容量执行与客流守恒}

路径搜索只产生移动意图，实际通过人数由共享容量执行层统一分配。对同一共享设施 \(r\)，无论移动请求来自哪条线路、哪个来源组或哪条边，只要它们映射到同一资源，就共同使用当前时间步的整数服务能力 \(C_r(t)\)。当请求人数超过服务能力时，未被服务的人数保留在上游队列。该处理保证同一设施不会因为被多条边引用而重复释放容量。

共享容量执行层还检查下游空间接收状态。若下游节点达到接收上限，即使设施本身仍有剩余服务能力，相关移动也不能直接释放到下游空间，而应作为空间阻塞或上游等待记录。执行完成后，系统更新节点人数、来源组分布、设施队列、在途批次和出口统计。由此，模型保持客流守恒：任一批次在同一时间步内只能处于上游等待、设施服务、边内移动、下游节点或出口中的一种状态。

共享容量执行层对 AA 和计算对照方法完全一致。这样设置使实验比较的差异限定在路径组织逻辑，即是否预测到达时刻设施负荷、是否写入未来到达事件、是否在选择阶段重评估路径，而不是来自设施容量、速度模型、空间接收规则或终止条件差异。

\subsection{4.4 计算对照方法}

ImprovedAStar 在本文中仅作为计算对照方法，不代表龙阳路站运营方制定或发布的官方疏散预案。该方法在同一车站网络、同一客流输入和同一共享容量执行层下进行路径搜索，但不使用未来设施到达事件集合，不预测候选批次到达设施时的排队状态，也不执行基于收益阈值的路径重评估。

两种方法采用共同终止规则和共同指标体系。仿真以目标客流全部到达出口为终止条件，记录整体疏散时间、尾部清空时间、停滞人秒、关键设施队列、来源组清空时间、出口利用和计算耗时。若 AA 缩短尾部时间但增加移动距离或计算时间，结果部分需要同时报告该权衡；若低负荷条件下两种方法表现接近，则该现象应作为方法适用边界的一部分说明。

\begin{table}[htbp]
\centering
\caption{AA 与计算对照方法的比较维度}
\small
\begin{tabular}{p{0.24\textwidth}p{0.30\textwidth}p{0.34\textwidth}}
\toprule
比较维度 & AA & ImprovedAStar \\
\midrule
路径评价状态 & 使用当前状态和已承诺未来到达事件 & 主要基于当前网络状态和边权评价 \\
设施等待处理 & 预测批次到达设施时的队列和等待 & 不显式预测到达时刻队列 \\
批次间相互影响 & 前序批次路径确定后写入未来事件，影响后续批次 & 同一轮批次之间不登记未来设施负荷 \\
执行层 & 与对照方法共用共享容量执行层 & 与 AA 共用共享容量执行层 \\
论文定位 & 本文提出的路径组织模型 & 计算基准，不是官方疏散方案 \\
\bottomrule
\end{tabular}
\end{table}

\section{5 龙阳路站案例研究与实验设计}

本文以龙阳路站疏散模型为案例验证 AA 的有效性。案例站包含 2 号线、7 号线、16 号线、18 号线和磁浮线之间的换乘关系，站内客流来源、换乘路径和共享设施类型较多，适合检验多来源客流在有限楼扶梯、通道、闸机和出口处汇聚时的路径组织效果。模型根据 CAD 图纸和仿真配置建立站台、站厅、换乘通道、楼扶梯、闸机和出口网络。设施宽度、闸机数量、出口宽度和服务能力在最终参数表中区分为实测值、CAD 或模型配置值、文献取值和假设值。

\subsection{5.1 案例车站建模与疏散场景}

案例建模包括空间网络构建、设施资源映射和客流来源设定三个部分。空间网络构建将站台、站厅、换乘通道、楼扶梯、闸机和出口抽象为节点与边；设施资源映射将共用楼扶梯、共用通道、闸机组和出口前区域登记为共享设施；客流来源设定将站台等待、站厅滞留、列车到达和换乘关系转化为自然到达批次。

实验设置低负荷常规应急场景和高负荷满载列车叠加场景。低负荷常规应急场景包含各线路站台等待、站厅和基础换乘客流，输入总人数为 2187 人。该场景用于检验设施竞争较弱时 AA 是否保持稳定，是否因过度预测造成不必要绕行或出口偏置。高负荷满载列车叠加场景在基础客流上加入多线路列车到达客流，输入总人数为 17905 人。该场景用于放大多线路站台客流、列车客流和换乘客流在共享设施上的汇聚效应，检验 AA 对等待暴露和尾部滞留的影响。

\begin{table}[htbp]
\centering
\caption{低负荷与高负荷疏散场景设置}
\small
\begin{tabular}{p{0.20\textwidth}p{0.16\textwidth}p{0.25\textwidth}p{0.27\textwidth}}
\toprule
场景 & 输入人数 & 主要客流构成 & 实验目的 \\
\midrule
低负荷常规应急 & 2187 & 站台等待、站厅滞留和基础换乘客流 & 检验设施竞争较弱时 AA 的稳定性和低负荷适用性 \\
高负荷满载列车叠加 & 17905 & 基础客流叠加多线路列车到达客流 & 检验多来源客流同步汇聚时 AA 对等待暴露和尾部滞留的影响 \\
\bottomrule
\end{tabular}
\end{table}

\subsection{5.2 到达负荷预测结果分析}

案例研究首先验证 AA 的到达负荷预测模块。该分析回答三个问题：第一，前序批次已经承诺的未来到达事件是否被计入后续批次的设施等待预测；第二，预计设施等待是否改变候选路径排序；第三，预测得到的关键设施压力是否能在执行层实际队列中得到对应体现。

实验选取换乘通道、楼扶梯、闸机和出口前区域中的关键共享设施作为分析对象。对每个设施记录预测到达人数、预测队列、实际队列、实际通过人数、未获得服务能力人数和最后排队时间。对代表性来源组批次，进一步分解候选路径代价，包括物理行走时间、设施等待时间、当前批次服务时间、空间回溢等待和总路径代价。该分析用于证明 AA 的中间机制是否真实参与路径决策，而不是只在最终结果中表现为一个黑箱差异。

\begin{table}[htbp]
\centering
\caption{AA 到达负荷预测分析指标}
\small
\begin{tabular}{p{0.20\textwidth}p{0.34\textwidth}p{0.34\textwidth}}
\toprule
分析对象 & 指标 & 作用 \\
\midrule
关键楼扶梯与换乘通道 & 预测到达人数、预测队列、实际队列、峰值等待、最后排队时间 & 判断多来源客流是否在垂向设施和换乘连接处集中汇聚 \\
闸机与出口前区域 & 通过人数、队列人秒、出口分配、来源组构成 & 判断出口方向选择是否造成局部服务压力 \\
代表性来源组批次 & 路径代价分解、候选路径排序、路径重评估记录 & 解释路径改变来自行走距离、设施等待还是空间回溢 \\
\bottomrule
\end{tabular}
\end{table}

\subsection{5.3 路径组织方案性能分析}

在到达负荷预测模块得到验证后，进一步比较 AA 与 ImprovedAStar 的路径组织效果。整体性能指标包括 \(T_{95}\)、\(T_{100}\)、平均疏散时间、累计停滞人秒、资源排队人秒、空间阻塞人秒、总移动距离和计算耗时。低负荷场景重点观察 AA 是否保持稳定；高负荷场景重点观察尾部清空、停滞等待和关键设施排队是否发生变化。

\begin{table}[htbp]
\centering
\caption{低负荷与高负荷场景下的整体疏散性能记录表}
\small
\begin{tabular}{p{0.18\textwidth}p{0.16\textwidth}p{0.10\textwidth}p{0.12\textwidth}p{0.12\textwidth}p{0.14\textwidth}p{0.12\textwidth}}
\toprule
场景 & 方法 & 人数 & \(T_{95}\) & \(T_{100}\) & 累计停滞人秒 & 总移动距离 \\
\midrule
低负荷常规应急 & ImprovedAStar & 2187 & 待冻结 & 待冻结 & 待冻结 & 待冻结 \\
低负荷常规应急 & AA & 2187 & 待冻结 & 待冻结 & 待冻结 & 待冻结 \\
高负荷满载列车叠加 & ImprovedAStar & 17905 & 待冻结 & 待冻结 & 待冻结 & 待冻结 \\
高负荷满载列车叠加 & AA & 17905 & 待冻结 & 待冻结 & 待冻结 & 待冻结 \\
\bottomrule
\end{tabular}
\end{table}

整体指标之后，结果分析转向关键设施瓶颈。若 AA 缩短 \(T_{100}\)，需要进一步说明这种变化是否来自关键楼扶梯、换乘通道、闸机或出口前区域的持续排队减少；若 AA 降低队列等待但增加移动距离或计算耗时，也需要同时报告该权衡。若低负荷场景下两种方法差异较小，则说明到达负荷预测主要在共享设施竞争明显时发挥作用。

\subsection{5.4 换乘客流汇聚机制分析}

换乘客流机制分析用于解释整体疏散结果的来源。对每个关键共享设施，统计站台等待、列车到达、站厅滞留和换乘来源的通过量、等待人秒和峰值队列贡献。若某一设施的队列主要由多个来源组叠加形成，说明该设施是全站疏散组织中的共享限制环节；若路径组织改变只是把等待从一个设施转移到另一个设施，则需要在设施链和来源组分解中呈现。

对换乘客流，进一步统计其来源线路、主要路径链、关键设施通过量、出口选择和清空时间。该分析回答三个问题：换乘客流是否是主要瓶颈形成的重要组成部分；AA 是否改变了换乘客流到达关键设施的时间分布；这种到达分布变化是否进一步影响全站整体疏散效率。通过这一层分析，换乘客流被作为解释全站疏散机制的核心变量，而不是被独立写成单一人群的疏散结果。

\subsection{5.5 Pathfinder 对照与适用边界}

Pathfinder 对照用于检验宏观网络模型的路径组织结果在微观行人仿真中是否呈现一致趋势。本文将宏观模型导出的来源组、路径链和路径人数转换为 Pathfinder 中的行人组和行为路径，并在相同车站几何和出口设置下运行微观仿真。对照指标包括累计疏散曲线、尾部清空时间、拥堵时间、关键设施附近空间拥挤分布和出口利用。

该对照不把 Pathfinder 视为真实车站观测，也不以 Pathfinder 结果替代宏观模型结果。其作用是检查宏观模型识别出的瓶颈位置是否在微观轨迹中出现相近拥堵，路径组织引起的出口和设施使用变化是否能在微观模型中复现。若两者存在差异，则从导航网格、个体避让、门选择规则和容量映射差异解释。

适用边界分析围绕三类变量展开：换乘客流比例、列车到达重叠程度和关键共享设施服务能力。换乘客流比例用于检验 AA 是否主要在换乘客流显著参与设施竞争时产生差异；列车到达重叠程度用于检验多线路客流同步释放对共享设施排队和尾部清空的影响；关键共享设施服务能力用于检验路径组织效果对瓶颈强度的依赖程度。

\begin{table}[htbp]
\centering
\caption{适用边界分析的变量设置}
\small
\begin{tabular}{p{0.22\textwidth}p{0.28\textwidth}p{0.38\textwidth}}
\toprule
变量类别 & 设置方式 & 分析目的 \\
\midrule
换乘客流比例 & 在总客流规模基本一致的前提下调整换乘来源占比 & 判断换乘客流汇聚是否是 AA 产生差异的重要场景条件 \\
列车到达重叠程度 & 调整多线路列车客流释放的时间间隔 & 判断同步到达和错峰到达对共享设施排队的影响 \\
关键共享设施服务能力 & 对主要楼扶梯、换乘通道或闸机服务率进行扰动 & 判断路径组织效果对瓶颈强度的依赖程度 \\
\bottomrule
\end{tabular}
\end{table}

最终结果冻结后，案例研究按以下顺序呈现：首先报告到达负荷预测结果，证明 AA 的中间预测机制确实改变路径评价；其次比较低负荷和高负荷场景下两种路径组织方法的整体疏散性能；再次通过关键设施和换乘客流来源分解解释整体差异的站内机制；最后通过 Pathfinder 对照和适用边界分析说明结果的跨模型一致性与适用范围。

\end{document}
